\section{Conclusiones de la etapa y trabajo pendiente}
\label{sec:conclusiones}

\subsection{Cumplimiento de los objetivos de la etapa}

\begin{longtable}{@{}L{5.0cm}C{1.6cm}L{6.9cm}@{}}
\filaencab \encab{Criterio de la pauta} & \encab{Peso} & \encab{Estado} \\ \midrule \endhead
Definición clara del problema & 10\% & Redactada con árbol de causas y efectos, justificación de relevancia y análisis de la oferta regional. \textbf{Pendiente el levantamiento directo con las partes interesadas} \\
Objetivos específicos y de relevancia & 25\% & Objetivo general y seis objetivos específicos con desglose SMART, indicador numérico, plazo y matriz de trazabilidad \\
Análisis de requerimientos & 45\% & 68 requerimientos con fuente y causa declaradas; siete alternativas de protocolo y cinco de implementación evaluadas mediante matrices ponderadas \\
Planificación de los objetivos & 20\% & 63 elementos con criterios de aceptación en el núcleo, estimaciones y verificación de viabilidad \\
\bottomrule
\caption{Estado de cumplimiento de los criterios de la Etapa 1.}
\label{tab:cumplimiento}
\end{longtable}

\subsection{Conclusiones principales}

\subsubsection{Sobre la elección del protocolo}

El análisis comparativo de siete alternativas arroja un resultado que no era evidente al inicio:
\textbf{ningún protocolo de telemetría contemporáneo combina compacidad extrema con vocabulario
semántico normalizado}. Los formatos compactos ---CayenneLPP--- estandarizan tipos de dato pero no
eventos ni estados de salud del dispositivo; los formatos semánticamente ricos ---LwM2M, Sparkplug B---
no caben en el payload disponible de LoRaWAN.

Contact ID resuelve ambas exigencias por una razón histórica: fue diseñado en la década de 1980 para
transmitirse mediante tonos telefónicos sobre una línea analógica, restricción que impuso un formato de
16 dígitos y un vocabulario cerrado de códigos. \textbf{Esa restricción de hace cuatro décadas es
precisamente la que lo hace apropiado hoy para radio de bajo consumo.}

\subsubsection{Sobre el vocabulario de salud de enlace}

Los códigos 381, 384, 344, 602 y 603 conforman un vocabulario normalizado de supervisión de nodos
inalámbricos alimentados a batería que precede al término \enquote{Internet de las Cosas}. La industria
de la seguridad electrónica resolvió y estandarizó ese problema décadas atrás. El aporte de este trabajo
consiste en reutilizar dicho vocabulario en un contexto abierto, en lugar de proponer uno nuevo que
incrementaría la fragmentación que se pretende resolver.

\subsubsection{Sobre la limitación identificada}

Contact ID no puede transportar valores analógicos. La limitación se declara y se resuelve mediante una
arquitectura de dos planos con requisitos de entrega diferenciados, decisión que además resulta
defendible en términos de ingeniería: los eventos exigen entrega confirmada y la telemetría admite
pérdida tolerable.

\subsubsection{Sobre la viabilidad del alcance}

La verificación numérica del backlog produjo el hallazgo de mayor consecuencia práctica de esta etapa:
\textbf{el alcance completo excede el tiempo disponible en un 37\%}. Detectarlo en la fase de
planificación, y no durante la implementación, permite acotar el compromiso al núcleo demostrativo, que
sí cabe con holgura. Esta constatación justifica retrospectivamente la decisión de estimar en horas y no
en unidades abstractas.

\subsection{Trabajo pendiente}

\begin{longtable}{@{}C{1.4cm}L{6.8cm}L{5.5cm}@{}}
\filaencab \encab{} & \encab{Pendiente} & \encab{Consecuencia si no se resuelve} \\ \midrule \endhead
\textbf{P1} & \textbf{Entrevistas a las partes interesadas} de la organización destinataria & Ausencia de validación directa del problema y del respaldo de los requerimientos \\
\textbf{P2} & \sout{Confirmación de las fechas de las cuatro etapas} \textbf{RESUELTO} & Calendario oficial publicado el 24/08: 4 sep, 2 oct, 30 oct, 4 dic \\
\textbf{P3} & Validación de la arquitectura de dos planos con el profesor guía & Rediseño del tratamiento de la telemetría \\
\textbf{P4} & Determinar si la normativa chilena impone restricción de tiempo de permanencia en 915--928\,MHz & Ninguna sobre la viabilidad; solo afecta el margen disponible \\
\textbf{P5} & \sout{Acceso al texto íntegro de los estándares SIA DC-05 y DC-09} \textbf{RESUELTO} & En agosto de 2025 la SIA publicó una implementación de referencia de DC-09 de código consultable (licencia \emph{Polyform Noncommercial}, que admite uso académico), lo que hace innecesario adquirir la norma \\
\textbf{P7} & Límite de longitud del campo de datos extendidos de DC-09 e identificadores ya asignados por convención & Podría obligar a acortar la codificación de valores o reasignar identificadores del perfil \\
\textbf{P8} & Gramática exacta de SIA DC-03 para valores numéricos, si se obtiene acceso al estándar & Si resultara más eficiente de lo estimado, DC-03 pasaría a ser la alternativa preferente \\
\textbf{P6} & Confirmar la disponibilidad de hardware de sensores y su presupuesto & Activación del riesgo R1; mitigado mediante emulación por software \\
\bottomrule
\caption{Trabajo pendiente al cierre de la Etapa 1.}
\label{tab:pendientes}
\end{longtable}

\subsection{Trabajo previsto para la Etapa 2}

Conforme a la fase \emph{Design} de PPDIOO y al backlog establecido, la Etapa 2 comprende:

\begin{itemize}
\item Diagramas de arquitectura de componentes, topología de red con segmentación, modelo de datos,
flujo de ingesta y secuencia de aprovisionamiento.
\item Especificación formal del perfil CID-IoT: gramática en notación ABNF, tabla de mapeo completa y
contrato de la interfaz de programación.
\item Determinación del modelo de despliegue y cuantificación del costo de operación (OE4).
\item Puesta en marcha del entorno de desarrollo e inicio de la implementación, adelantando trabajo
respecto de la Etapa 3 conforme a la redistribución correctiva de la sección~\ref{sec:viabilidad}.
\end{itemize}
